% Options for packages loaded elsewhere
\PassOptionsToPackage{unicode}{hyperref}
\PassOptionsToPackage{hyphens}{url}
%
\documentclass[
]{article}
\usepackage{amsmath,amssymb}
\usepackage{iftex}
\ifPDFTeX
  \usepackage[T1]{fontenc}
  \usepackage[utf8]{inputenc}
  \usepackage{textcomp} % provide euro and other symbols
\else % if luatex or xetex
  \usepackage{unicode-math} % this also loads fontspec
  \defaultfontfeatures{Scale=MatchLowercase}
  \defaultfontfeatures[\rmfamily]{Ligatures=TeX,Scale=1}
\fi
\usepackage{lmodern}
\ifPDFTeX\else
  % xetex/luatex font selection
\fi
% Use upquote if available, for straight quotes in verbatim environments
\IfFileExists{upquote.sty}{\usepackage{upquote}}{}
\IfFileExists{microtype.sty}{% use microtype if available
  \usepackage[]{microtype}
  \UseMicrotypeSet[protrusion]{basicmath} % disable protrusion for tt fonts
}{}
\makeatletter
\@ifundefined{KOMAClassName}{% if non-KOMA class
  \IfFileExists{parskip.sty}{%
    \usepackage{parskip}
  }{% else
    \setlength{\parindent}{0pt}
    \setlength{\parskip}{6pt plus 2pt minus 1pt}}
}{% if KOMA class
  \KOMAoptions{parskip=half}}
\makeatother
\usepackage{xcolor}
\setlength{\emergencystretch}{3em} % prevent overfull lines
\providecommand{\tightlist}{%
  \setlength{\itemsep}{0pt}\setlength{\parskip}{0pt}}
\setcounter{secnumdepth}{-\maxdimen} % remove section numbering
\ifLuaTeX
  \usepackage{selnolig}  % disable illegal ligatures
\fi
\usepackage{bookmark}
\IfFileExists{xurl.sty}{\usepackage{xurl}}{} % add URL line breaks if available
\urlstyle{same}
\hypersetup{
  pdftitle={Tumor-informed versus tumor-agnostic approaches in addressing minimal residual disease --- unraveling the optimal strategy (writing)},
  hidelinks,
  pdfcreator={LaTeX via pandoc}}

\title{Tumor-informed versus tumor-agnostic approaches in addressing
minimal residual disease --- unraveling the optimal strategy (writing)}
\author{}
\date{}

\begin{document}
\maketitle

\subsection{II. Tumor-agnostic
strategies}\label{ii.-tumor-agnostic-strategies}

Purely methylation-based tests have sucessfuly been used to assess MRD
in a tumor-agnostic manner. Indeed, these biomarkers hold great value
for this specific approach, because of their widespread and common
occurrence in cancer, which contrasts the higher tumor and clonal
specificity that genetic variants exhibit
{[}@haoDNAMethylationMarkers2017; @luoCirculatingTumorDNA2020b{]}.

\subsubsection{WIF1 and NPY}\label{wif1-and-npy}

One of the most relevant examples are the WIF1 and NPY genes, which have
consistently been found to be hypermethylated in CRC tissue and have
been proposed as blood biomarkers for ctDNA and MRD detection
{[}@roperchAberrantMethylationNPY2013;
@garrigouStudyHypermethylatedCirculating2016{]}. For instance, one study
assessed the methylation status of these genes via digital PCR in a
retrospective cohort of 1345 stage III CRC patients from the IDEA-FRANCE
clinical trial (NCT00958737). A total of 140 patients tested positive
for ctDNA (13.8\%), and their 3-year DFS proved to be significantly
lower than the ctDNA-negative group (66.39\% vs. 76.71\%)
{[}@taiebPrognosticValueRelation2021{]}. A digital PCR platform was also
able to achieve 95.5\% sensitivity and 100\% specificity in a cohort of
22 liquid biopsies {[}@oversDetectionSpecificHypermethylated2021a{]}.

\subsubsection{BCAT1 and IKZF1}\label{bcat1-and-ikzf1}

The methylation levels of BCAT1 and IKZF1 in ctDNA has also been
correlated to CRC tumor presence. In a prospective study of 142 stage
I-III CRC patients, the aberrant methylation of these two genes was
found in 19 patients (13.38\%), which exhibited higher risk of
recurrence (HR = 5.7) and lesser 3-year RFS (56.5\% vs. 83.3\% in
ctDNA-negative patients) {[}@pedersenDetectionMethylatedBCAT12022{]}.
The relevance of BCAT1 and IKZF1 has been leveraged by the
laboratory-developed test Colvera (CSIRO), which in its latest iteration
has exhibited values of sensitivity and specificity as high as 62\% and
98.3\%, respectively
{[}@vuceticClinicalPerformanceMethylationbased2021{]}.

\subsubsection{SEPT9}\label{sept9}

SEPT9, which has extensively been studied for CRC early detection and is
at the base of the FDA-approved Epi proColon commercial assay
(Epigenomics AG) has also been proposed for MRD detection
{[}@wangAdvancePlasmaSEPT92018{]}. After tumoral resection, the
detection of methylated SEPT9 was able to predict the risk of recurrence
in a cohort of 426 stage II-II CRC patients, among which RFS was
decreased (HR = 5.83) {[}@yuanCirculatingMethylatedSEPT92022{]}.

\subsubsection{Genome-wide}\label{genome-wide}

\href{https://www.ncbi.nlm.nih.gov/pmc/articles/PMC10541738/}{Response
prediction by mutation- or methylation-specific detection of ctDNA
dynamics in pretreated metastatic colorectal cancer - PMC}

\subsection{III. Tumor-informed
strategies}\label{iii.-tumor-informed-strategies}

As previously stated, the tumor-informed approach, aims to molecularly
characterize the tumor tissue prior to ctDNA detection. This way, a
bespoke ctDNA assay can be created to directly target the specific
mutations of said tumor. The addition of tumor sequencing has both
advantages and disadvantages when compared to the tumor-agnostic
strategy.

{[}CHIP{]}

However, the addition of tumor sequencing comes at a cost, both in terms
of he

Focusing only on variants found in tumor tissue also raises the question
of tumor heterogeneity and resistance mechanisms, where treatment may
lead to the appearance or increase of different clonal cells harboring
different mutations than those initially found. Thus, potential new
variants could be omitted.

\subsubsection{DYNAMIC trial (SafeSeqS
approach)}\label{dynamic-trial-safeseqs-approach}

\begin{quote}
Añadir sobre la capacidad de detección.\\
El primero de TIE van solo a por una y el segundo por al menos una???
revisar
\end{quote}

A growing body of research and clinical trial results has revealed the
usefulness of tumor-informed approaches to MRD detection.

One of the earliest key studies that supported this approach was
performed in a prospective cohort of 230 CRC stage II patients (trial
ACTRN12612000326897), where ctDNA detection after resection showed
higher risk of recurrence as determined by a CT scan. Initially, the
variant with the highest allele frequency in FFPE tumor tissue was
determined for each patient using targeted multiplex PCR aiming for 15
genes relevant to CRC cancer. Afterwards, those variants were tracked
using a personalized Safe-SeqS assay, an UMI-based approach that enables
deep sequencing with error correction
{[}@kindeDetectionQuantificationRare2011{]}. Patients were then
classified as ctDNA-positive or -negative, depending on the finding of
said mutations in plasma. The results showed that postoperative
ctDNA-positive status was related to higher risk of recurrence among
patients not treated with adjuvant chemotherapy (HR = 18; 95\%CI =
7.9-40) {[}@tieCirculatingTumorDNA2016{]}.

The same tumor-informed approach was adopted in the DYNAMIC trial
(ACTRN12615000381583), which further studied the possibility of using
the postoperative ctDNA status as a basis for making treatment decisions
for patients with CRC stage II. By random assignment, patients were
assigned to an standard management group (n = 153) or a ctDNA-guided
management group (n = 302). In this approximation, patients that were
ctDNA-positive at 4 or 7 weeks after surgery were recommended to receive
adjuvant oxaliplatin- or fluoropyrimidine-based chemotherapy.
Conversely, patients testing negative for ctDNA did not receive adjuvant
chemotherapy. The RFS for the ctDNA-guided group was comparable to that
of the standard management group (93.5\% vs. 92.4\%), indicating the
suitability of chemotherapy reduction in ctDNA negative patients
{[}@tieCirculatingTumorDNA2022{]}.

\subsubsection{Signatera assay: GALAXY trial and BESPOKE
trial}\label{signatera-assay-galaxy-trial-and-bespoke-trial}

The Signatera assay (Natera) has become a prominent advocate of the
tumor-informed approach. In brief, whole exome sequencing is performed
on both primary tumor tissue and matched normal tissue to work out a set
of 16 or more mutations. These are subsequently considered for the
creation of a bespoke test based on ultradeep multiplex PCR, which is
used to detect ctDNA and monitorize MRD.

One key study that employs this approach in CRC is the GALAXY trial
(UMIN000039205), part of the CIRCULATE-Japan study. Here, a cohort of
1039 patients with resectable CRC (stages II-IV) was subject to an
observational study, with postoperative ctDNA status determination and
follow-up to a median of 16.74 months. After surgery, the ctDNA status
was used as the criterion for patient assignment to one of two related
clinical trials: patients that tested positive for ctDNA were assigned
to the ALTAIR trial, assessing the adequacy of chemotherapy escalation,
and ctDNA-negative patients were assigned to the VEGA trial, studying
the noninferiority of treatment deescalation. The ctDNA determination
exhibited prognostic value, with patients that tested positive at 4
weeks after surgery showing an increased risk of recurrence (HR = 10.0)
{[}@kotaniMolecularResidualDisease2023{]}.

The Signatera assay is also the centerpiece of the ongoing BESPOKE
clinical trial (NCT04264702).{[}@kasiBESPOKEStudyProtocol2021{]}.

\begin{quote}
thereby expanding upon the results of the aforementioned DYNAMIC trial.
\end{quote}

\href{https://dailynews.ascopubs.org/do/studies-confirm-utility-ctdna-guiding-adjuvant-chemotherapy-crc}{Studies
Confirm the Utility of ctDNA in Guiding Adjuvant Chemotherapy in CRC
\textbar{} ASCO Daily News}\^{}\^{}\^{}\^{}\^{} First results of the
Signatera assay

\subsubsection{Other trials}\label{other-trials}

Studies using non-commercial tumor-informed assays have also rendered
insightful results. One study used a custom targeted panel of 29 genes
to study tumor tissue mutations, and performed a ctDNA-based follow-up
of said variants using ddPCR. The detection of ctDNA after surgery was
significantly associated to lower disease free survival (HR = 6.96), and
so was the tracking across multiple postsurgical samples (HR = 8.03)
{[}@tarazonaTargetedNextgenerationSequencing2019{]}.

Other recent studies have explored the

\subsection{IV. Comparison}\label{iv.-comparison}

\subsubsection{Strengths and weaknesses}\label{strengths-and-weaknesses}

High specificity of mutations may lead to false negatives in
tumor-agnostic assays. However,

By sequencing both the tumor tissue and the ctDNA, variants related to
clonal hematopoiesis of indeterminate potential can be directly filtered
{[}@steensmaClonalHematopoiesisIndeterminate2015{]}.

\subsubsection{Discussion of studies comparing both
approaches}\label{discussion-of-studies-comparing-both-approaches}

{Tumor-informed versus tumor-agnostic approaches in addressing minimal
residual disease --- unraveling the optimal strategy (planning)
\textgreater{} Studies comparing both approaches}

\subsection{V. Emerging Technologies and Future
Directions}\label{v.-emerging-technologies-and-future-directions}

\begin{itemize}
\tightlist
\item
  Fragmentomics
\item
  CTCs
\item
  Imaging techniques (?) (guided by AI?)
\end{itemize}

\subsubsection{1. Genomics, epigenomics and
fragmentomics}\label{genomics-epigenomics-and-fragmentomics}

The study of whole genome methylation markers through techniques such as
cfMeDIP-seq ---which enables high resolution profiling at a relatively
low cost when compared with bisulfite methods--- is also starting to
render insightful results.

The study of cfDNA fragmentation patterns ---a field termed
\textbf{fragmentomics}--- is also starting to gain traction in the
setting of tumor-agnostic MRD detection.

A study in 51 metastatic cancer patients from the INSPIRE trial
(NCT02644369) that followed a tumor-agnostic approach showed that cancer
specific methylation predicted OS and PFS similarly to mutation
concentration in ctDNA. When integrating fragmentomics, both OS and PFS
prediction surpassed that of mutation concentration
{[}@zhao1664MOTumornaiveMethylomes2022{]}.

The integration of multiple omics is certainly a target to pursue for
the promise of a more sensitive and specific detection of ctDNA.
\href{https://doi.org/10.1016/j.annonc.2022.07.1744}{Redirecting}.
Similar approaches could be followed for MRD detection, although the
high related costs would be currrently unaffordable for the daily
clinical practice.

EC:

\begin{itemize}
\tightlist
\item
  ECLIPSE
\item
  GRAIL ??\textquotesingle{}
\item
  MEDAL
\end{itemize}

The potential of

\subsubsection{2. Liquid Biopsy Technologies:
CTCs}\label{liquid-biopsy-technologies-ctcs}

\subsubsection{3. Advanced Imaging
Techniques}\label{advanced-imaging-techniques}

\begin{quote}
The application of AI in radiomics.\\
However, ctDNA still takes
\end{quote}

{[}@tibermacineRadiomicsModellingRectal2021{]}

\subsubsection{4. AI}\label{ai}

\begin{quote}
Beyond radiomics, AI also has interesting applications...
\end{quote}

\end{document}
